\documentclass{article}
\usepackage{graphicx} % Required for inserting images
\usepackage{amsmath}
\usepackage{algorithmic}
\usepackage{algorithm}
\usepackage[polish]{babel}
\usepackage[utf8]{inputenc}
\usepackage[T1]{fontenc}
\usepackage{float}
\usepackage{listings}
\usepackage[table]{xcolor}
\usepackage{diagbox}
\usepackage{geometry}

\title{Algorytmy Równoległe - Laboratorium 4}
\author{Adam Staniszewski}

\begin{document}

\maketitle

\section*{Przygotowanie danych}
Wykonajmy pomiary czasu wykonania algorytmu dla implementacji sekwencyjnej (czyli dla jednego rdzenia) i równoległej (dwa i cztery rdzenie). Wykorzystamy uśredniony czas wykonania z pięciu prób, Rozważymy też różne rozmiary problemów.
\begin{table}[ht]
\centering
\begin{tabular}{|l|r|r|r|}
\hline
\rowcolor{orange!40} 
Liczba rdzeni & Średni czas wykonania algorytmu (s) \\\hline
1 (implementacja sekwencyjna) & 0.018625 \\\hline
2 & 0.032318 \\\hline
4 &  0.026899 \\\hline

\end{tabular}
\caption{\label{tab:widgets}Średni czas obliczeń dla problemu o rozmiarze 10 na 10.}
\end{table}

\begin{table}[ht]
\centering
\begin{tabular}{|l|r|r|r|}
\hline
\rowcolor{orange!40} 
Liczba rdzeni & Średni czas wykonania algorytmu (s)\\\hline
1 (implementacja sekwencyjna) & 14.41022 \\\hline
2 & 8.472051 \\\hline
4 & 3.904285 \\\hline

\end{tabular}
\caption{\label{tab:widgets}Średni czas obliczeń dla problemu o rozmiarze 50 na 50.}
\end{table}

\begin{table}[ht]
\centering
\begin{tabular}{|l|r|r|r|}
\hline
\rowcolor{orange!40} 
Liczba rdzeni & Średni czas wykonania algorytmu (s)\\\hline
1 (implementacja sekwencyjna) & 207.95182 \\\hline
2 & 115.18612 \\\hline
4 & 52.76849 \\\hline

\end{tabular}
\caption{\label{tab:widgets}Średni czas obliczeń dla problemu o rozmiarze 100 na 100.}
\end{table}

\section*{Wykresy}
\subsection*{Przypieszenie}

\begin{figure}[H]
    \centering
    \includegraphics[width=1\linewidth]{speedup.png}
    \caption{Wartości przyspieszenia dla różnych rozmiarów problemu i różnej liczby rdzeni}
    \label{fig:enter-label}
\end{figure}

\subsection*{Efektywność}

\begin{figure}[H]
    \centering
    \includegraphics[width=1\linewidth]{effectivness.png}
    \caption{Wartości efektywności dla różnych rozmiarów problemu i różnej liczby rdzeni}
    \label{fig:enter-label}
\end{figure}

\subsection*{Metryka Karpa-Flatta}

\begin{figure}[H]
    \centering
    \includegraphics[width=1\linewidth]{karp-flatt.png}
    \caption{Wartości metryki Karpa-Flatta dla różnych rozmiarów problemu i różnej liczby rdzeni}
    \label{fig:enter-label}
\end{figure}

\subsection*{Wnioski}
Możemy zauważyć, że dla małych problemów (rozmiaru 10x10) wartości przyspieszenia i efektywności spadają wraz z liczbą rdzeni. Możemy zatem wywnioskować, że dla dostatecznie małych problemów czas wykonania algorytmu jest zdominowany przez czas komunikacji między procesami. Ujemne wartości metryki Karpa-Flatta dla problemu o rozmiarze 10x10 potwierdzają powyższe przypuszczenie.

Dla większych problemów możemy zauważyć, że wszystkie metryki są rosnące. Dodatnie wartości metryki Karpa-Flatta wskazują, że zrównoleglenie algorytmu jest korzystne dla większych problemów.

\end{document}
